% 第1回は ../class1/class1PT26.tex、第2・3回は ../class2/class2PT26.tex、
% 第4回は class4PT26.tex のレポート問題を収録する。
% 第2回は resume2PT26.tex では省略されている第3・4問も収録する。
% Build in this directory (run twice for page references):
% pdflatex -interaction=nonstopmode -halt-on-error -file-line-error report.tex
\documentclass[a4paper,10pt,pdflatex,ja=standard,jafont=ipaex-type1,
  magstyle=real,scale=1]{bxjsarticle}
\usepackage{lmodern}
\usepackage{amsmath,amssymb,mathtools}
\usepackage{geometry}
\geometry{reset,left=22mm,right=22mm,top=20mm,bottom=17mm,
  headheight=13pt,headsep=6mm,footskip=10mm}
\usepackage{xcolor}
\usepackage{enumitem}
\usepackage{array,booktabs,tabularx}
\usepackage{fancyhdr}
\usepackage{hyperref}

\definecolor{ink}{HTML}{193A5C}
\definecolor{accent}{HTML}{237D82}
\definecolor{soft}{HTML}{EAF3F3}
\definecolor{muted}{HTML}{52616B}
\newcommand{\coursename}{確率論と統計学 I}
\hypersetup{unicode=true,colorlinks=true,urlcolor=accent,linkcolor=ink,
  pdftitle={第1〜4回 レポート問題と解答}}
\pagestyle{fancy}
\fancyhf{}
\fancyhead[L]{\small\color{ink}\coursename{}・レポート問題と解答}
\fancyhead[R]{\small\color{ink}第4回配布（第1〜4回）}
\fancyfoot[C]{\small\color{ink}\thepage/\pageref*{lastpage}}
\renewcommand{\headrulewidth}{0.3pt}
\setlength{\parindent}{0pt}
\setlength{\parskip}{2pt}
\setlength{\abovedisplayskip}{4pt}
\setlength{\belowdisplayskip}{4pt}
\setlength{\abovedisplayshortskip}{3pt}
\setlength{\belowdisplayshortskip}{4pt}
\setlist[itemize]{leftmargin=1.8em,itemsep=1pt,topsep=2pt,parsep=0pt}
\setlist[enumerate]{leftmargin=2em,itemsep=3pt,topsep=3pt,parsep=0pt}
\newcommand{\topic}[1]{\par\vspace{5pt}
  {\large\sffamily\mdseries\color{accent}#1}\par\vspace{2pt}}
\newcommand{\handouttitle}[1]{%
  \begin{center}
    {\Large\sffamily\mdseries\color{ink}#1}\par\vspace{3mm}
    {\color{accent}\rule{0.88\linewidth}{0.7pt}}
  \end{center}\vspace{-2mm}}
\newcommand{\keybox}[1]{\par\vspace{3pt}\begingroup
  \setlength{\fboxsep}{5pt}
  \colorbox{soft}{\parbox{\dimexpr\linewidth-2\fboxsep\relax}{#1}}%
  \endgroup\par\vspace{2pt}}
\newcommand{\examplelabel}[1]{{\sffamily\mdseries\color{ink}#1}}
\newcommand{\F}{\mathcal F}
\newcommand{\R}{\mathbb R}
\newcommand{\Prob}{\mathbb P}
\newcommand{\E}{\mathbb E}
\newcommand{\ind}{\mathrm{1}}

\begin{document}
\fontsize{10pt}{13.5pt}\selectfont
\handouttitle{第1回\quad 確率空間：レポート問題と解答}
{\small\color{muted}第1〜4回の全12問を収録する。問題番号「$i.j$」は第$i$回の第$j$問を表す。}

\topic{問題1.1\quad 天気の確率空間}
\keybox{晴れの確率が30\%、くもりが50\%、雨が20\%となる確率空間を構成せよ。
（ただし$\mathbb P$を具体的に構成する必要はなく、$p_{\omega}$のみでよい）}

\examplelabel{解答}
標本空間、事象の集合族、各結果の重みを
$$
  \begin{gathered}
    \Omega=\{\text{晴れ},\text{くもり},\text{雨}\},\qquad \F=2^\Omega,\\
    p_{\text{晴れ}}=\frac3{10},\qquad
    p_{\text{くもり}}=\frac5{10},\qquad
    p_{\text{雨}}=\frac2{10}
  \end{gathered}
$$
と定める。$2^\Omega$ は $\Omega$ のすべての部分集合からなる集合族である。
各重みは非負で、
$$
  \sum_{\omega\in\Omega}p_\omega=\frac3{10}+\frac5{10}+\frac2{10}=1
$$
を満たすので、これらは離散確率空間を定める。具体的には
$\Prob(A)=\sum_{\omega\in A}p_\omega$ $(A\subseteq\Omega)$ とすればよい。

\topic{問題1.2\quad 確率を測る集合を制限する理由}
\keybox{一般の確率空間で、確率を測れる集合を $\F$ に制限する理由を説明せよ。}

\examplelabel{解答}
確率を測る事象については、補事象や可算個の和事象も扱いたい。
そのため、確率測度の定義域 $\F$ には
$$
  \Omega\in\F,\qquad
  A\in\F\ \Longrightarrow\ A^c\in\F,\qquad
  A_n\in\F\ (n\ge1)\ \Longrightarrow\ \bigcup_{n=1}^{\infty}A_n\in\F
$$
を要求する。これが $\sigma$-加法族の条件である。

有限・可算な標本空間では、各点の重みからすべての部分集合の確率を定められる。
一方、区間上の一様な確率を考える場合には、すべての部分集合に確率を与えようとすると、
可算加法性と一様性に対応する移動不変性が両立しない。問題1.3でその例を示す。
したがって、ボレル集合族などの適切な $\F$ を選び、その上で確率を定義する。

また、$\F$ は「どの事象を区別できるか」という情報の範囲を表すこともできる。
例えば $\F=\{\varnothing,\Omega\}$ は、結果の違いを区別しない情報に対応する。

\keybox{元の問題の注意：レポート問題の第2,3問はAIと相談してもよいが、その場合は
どのような議論をしたかも書くこと。ただし、内容を自分で理解したうえで、
自分の言葉で説明すること。}

\clearpage
\handouttitle{第1回\quad 確率空間：レポート問題と解答（続き）}

\topic{問題1.3\quad すべての部分集合で測れない例}
\keybox{第2問と関連して、自然に求める性質を保つために $\F$をすべての$\Omega$の部分集合
とできない例を一つ挙げ、どの性質が両立しないのか述べよ。}

\examplelabel{解答}
区間 $\Omega=[0,1)$ の両端をつないだ円周上の一様な確率を考える。
$t\in[0,1)$ に対して、円周に沿った移動を
$$
  T_t(x)=x+t-\lfloor x+t\rfloor\qquad(x\in[0,1))
$$
と定める。$\lfloor u\rfloor$ は $u$ 以下の最大の整数である。
一様性に対応する自然な条件は、事象を移動しても確率が変わらないこと、すなわち
$$
  \Prob(T_t(A))=\Prob(A),\qquad
  T_t(A):=\{T_t(x):x\in A\}
$$
である。ボレル集合上のルベーグ測度はこの条件を満たす。
以下では、この条件と可算加法性を保ったまま $\F=2^\Omega$ とすることはできないと示す。

$[0,1)$ 上に同値関係
$$
  x\sim y\quad\Longleftrightarrow\quad x-y\in\mathbb Q
$$
を定める。選択公理を用いて各同値類から一点ずつ選び、その集合を $V$ とする。
これはヴィタリ集合の円周上での構成である。
$Q=\mathbb Q\cap[0,1)$ とおき、$q\in Q$ に対して $V_q=T_q(V)$ と定める。

任意の $x\in[0,1)$ について、その同値類の代表元 $v\in V$ が存在する。
$x-v\in\mathbb Q$ なので、その差を1を法として $[0,1)$ に移した数 $q$ は $Q$ に属し、
$x=T_q(v)$ となる。したがって
$$
  [0,1)=\bigcup_{q\in Q}V_q.
$$
また、$T_q(v)=T_r(w)$ $(v,w\in V,\ q,r\in Q)$ なら $v-w\in\mathbb Q$ である。
代表元の選び方から $v=w$ となり、続いて $q-r\in\mathbb Z$ が従う。
$q,r\in[0,1)$ より $q=r$ である。よって、$V_q$ は互いに素である。

もしすべての部分集合に確率を与えられるなら、$V$ とすべての $V_q$ も事象となる。
移動不変性より $\Prob(V_q)=\Prob(V)=:c\ge0$ であり、$Q$ は可算無限集合なので、
可算加法性から
$$
  1=\Prob([0,1))=\sum_{q\in Q}\Prob(V_q)=\sum_{q\in Q}c
$$
となる。しかし、右辺は $c=0$ なら $0$、$c>0$ なら $+\infty$ であり、矛盾する。

以上より、すべての部分集合への定義、非負性と全体の確率が1であること、
可算加法性、円周上の移動不変性は同時には満たせない。
なお、任意の集合 $\Omega$ 上でも、一点 $\omega_0$ に確率を集中させる
$$
  \Prob(A)=\ind_A(\omega_0)\qquad(A\subseteq\Omega)
$$
という確率測度は定義できる。問題となるのは、上で要求した一様性も保つ場合である。

\clearpage
\handouttitle{第2回\quad 確率変数と分布：レポート問題と解答}

\topic{問題2.1\quad コイン3回の裏の数}
\keybox{公平なコインを３回投げる。裏の数を表す確率変数を定義したい。
確率空間およびその確率変数の定義を述べよ。}

\examplelabel{解答}
各回のコイン投げは独立とする。$H$ を表、$T$ を裏として
$$
  \Omega=\{H,T\}^3,\qquad \F=2^\Omega,\qquad
  \Prob(\{\omega\})=\frac18\quad(\omega\in\Omega)
$$
と定める。すなわち、$\Prob(A)=|A|/8$ $(A\subseteq\Omega)$ である。
ここで $|A|$ は $A$ の要素数を表す。裏の数を表す写像を
$$
  X:\Omega\longrightarrow\R,\qquad
  X(\omega_1,\omega_2,\omega_3)=\sum_{j=1}^{3}\ind_{\{T\}}(\omega_j)
$$
と定義する。$\F=2^\Omega$ なので、任意のボレル集合 $B\subseteq\R$ について
$X^{-1}(B)\in\F$ となり、$X$ は確率変数である。
例えば $X(H,T,T)=2$、$X(H,H,H)=0$ である。
確認として、その分布は
$$
  \Prob(X=k)=\frac{\binom3k}{8}\quad(k=0,1,2,3),\qquad
  \Prob(X=k)=0\quad(k\notin\{0,1,2,3\})
$$
となる。公平であるという条件だけでは回ごとの独立性は定まらないため、
上では独立なコイン投げのモデルを明示した。

\topic{問題2.2\quad 逆像で定まる集合族}
\keybox{確率空間$(\Omega,\mathcal F,\mathbb P)$と写像 $X:\Omega\to\mathbb R$を考える。
次の集合族が$\sigma$加法族であることを示せ：
$$\{B\subset\mathbb R:\ X^{-1}(B)\in\mathcal F\}.$$}

\examplelabel{解答}
$\mathcal G:=\{B\subseteq\R:X^{-1}(B)\in\F\}$ とおく。
まず $X^{-1}(\R)=\Omega\in\F$ なので、$\R\in\mathcal G$ である。
$B\in\mathcal G$ に対して、補集合をそれぞれ $\R$ と $\Omega$ の中で取ると
$$
  X^{-1}(\R\setminus B)=\Omega\setminus X^{-1}(B)\in\F
$$
であり、$\R\setminus B\in\mathcal G$ となる。
さらに、$B_n\in\mathcal G$ $(n\ge1)$ なら
$$
  X^{-1}\!\left(\bigcup_{n=1}^{\infty}B_n\right)
  =\bigcup_{n=1}^{\infty}X^{-1}(B_n)\in\F,
  \qquad \bigcup_{n=1}^{\infty}B_n\in\mathcal G.
$$
以上より、$\mathcal G$ は $\R$ 上の $\sigma$-加法族である。
この証明では、$X$ があらかじめ確率変数であるとは仮定していない。

\clearpage
\handouttitle{第2回\quad 確率変数と分布：レポート問題と解答（続き）}

\topic{問題2.3\quad 自分で構成する確率変数と分布関数}
\keybox{離散確率空間、連続確率空間それぞれで確率変数を自分で構成し、
分布関数を求めよ。}

\examplelabel{解答例：離散確率空間}
公平なサイコロを1回投げるモデル
$$
  \Omega_d=\{1,2,3,4,5,6\},\qquad \F_d=2^{\Omega_d},\qquad
  \Prob_d(A)=\frac{|A|}{6}
$$
上で、偶数の目の指示関数 $X_d(\omega)=\ind_{\{2,4,6\}}(\omega)$ を定める。
すべての部分集合が事象なので、$X_d$ は確率変数であり、
$\Prob_d(X_d=0)=\Prob_d(X_d=1)=1/2$ である。したがって
$$
  F_{X_d}(x)=\Prob_d(X_d\le x)=
  \begin{cases}
    0,&x<0,\\
    \dfrac12,&0\le x<1,\\
    1,&x\ge1.
  \end{cases}
$$

\examplelabel{解答例：連続確率空間}
$\lambda$ をルベーグ測度として
$$
  \Omega_c=[0,1],\qquad \F_c=\mathcal B([0,1]),\qquad
  \Prob_c(A)=\lambda(A),\qquad X_c(\omega)=\omega^2
$$
と定める。$X_c$ は連続関数なのでボレル可測、すなわち確率変数である。
$0\le x<1$ のとき
$$
  \{\omega\in[0,1]:X_c(\omega)\le x\}=[0,\sqrt{x}]
$$
より、分布関数はすべての $x\in\R$ に対して
$$
  F_{X_c}(x)=
  \begin{cases}
    0,&x<0,\\
    \sqrt{x},&0\le x<1,\\
    1,&x\ge1
  \end{cases}
$$
となる。

\topic{問題2.4\quad 確率変数の和の可測性}
\keybox{同じ確率空間で定義された確率変数 $X,Y$に対し、
$X+Y$も確率変数であることを示せ。}

\examplelabel{解答}
$X,Y$ は実数値確率変数とする。任意の $a\in\R$ について
$$
  \{X+Y<a\}
  =\bigcup_{q\in\mathbb Q}\bigl(\{X<q\}\cap\{Y<a-q\}\bigr)
$$
が成り立つ。実際、$X(\omega)+Y(\omega)<a$ なら、
$X(\omega)<q<a-Y(\omega)$ を満たす有理数 $q$ が存在する。逆向きは不等式の加算による。
各集合 $\{X<q\},\{Y<a-q\}$ は $\F$ に属し、$\mathbb Q$ は可算なので
$\{X+Y<a\}\in\F$ である。半直線 $(-\infty,a)$ はボレル集合族を生成するから、
$X+Y$ は確率変数である。この結論に $X,Y$ の独立性は必要ない。

\clearpage
\handouttitle{第3回\quad 期待値と積分：レポート問題と解答}
以下の2問では、$\Omega=(0,1)$、$\F=\mathcal B((0,1))$、$P=\lambda$ とする。
$\lambda$ はルベーグ測度であり、この区間の長さは1である。

\topic{問題3.1\quad 積分による期待値の計算}
\keybox{$X(\omega)=\omega^2$ とする。$\E[X]$ を積分により求めよ。}

\examplelabel{解答}
$X$ は非負で有界な連続関数なので、可積分な確率変数である。
$$
  \E[X]=\int_{\Omega}X\,dP=\int_0^1\omega^2\,d\omega
  =\left[\frac{\omega^3}{3}\right]_0^1=\frac13.
$$
端点 $0,1$ はルベーグ測度0の集合なので、開区間 $(0,1)$ 上の積分を上の形で計算できる。

\topic{問題3.2\quad 点ごとの収束と優収束定理}
\keybox{$X_n(\omega)=n\ind_{(0,1/n)}(\omega)$ とする。
$X_n\to0$ と $\E[X_n]=1$ を示し、優収束定理を適用できない理由を述べよ。
（理由の部分はAIを用いてよいが、自分の言葉で書くこと）}

\examplelabel{解答}
$n$ は正の整数とする。各 $X_n$ はボレル集合の指示関数の定数倍なので確率変数である。
任意に $\omega\in(0,1)$ を固定し、$N>1/\omega$ となる整数 $N$ を選ぶ。
$n\ge N$ なら $1/n<\omega$ なので
$$
  X_n(\omega)=0\quad(n\ge N),\qquad
  \lim_{n\to\infty}X_n(\omega)=0.
$$
したがって、すべての $\omega\in\Omega$ で $X_n(\omega)\to0$ となる。
一方、各 $n\ge1$ について
$$
  \E[X_n]=\int_0^1 n\ind_{(0,1/n)}(\omega)\,d\omega
  =n\lambda((0,1/n))=n\cdot\frac1n=1.
$$
よって、$\lim_{n\to\infty}\E[X_n]=1\ne0=\E[\lim_{n\to\infty}X_n]$ である。

優収束定理には、すべての $n$ に共通する可積分な支配関数が必要である。
そのような非負確率変数 $G$ が存在し、$X_n\le G$ が各 $n$ についてほとんど確実に成り立つと仮定する。
互いに素な区間
$$
  I_m=\left(\frac1{m+1},\frac1m\right)\qquad(m=1,2,\ldots)
$$
上では $X_m=m$ なので、$G\ge m$ が $I_m$ 上でほとんど至るところ成り立つ。
したがって、任意の正の整数 $M$ に対して
$$
  \E[G]\ge\sum_{m=1}^{M}\int_{I_m}G\,d\lambda
  \ge\sum_{m=1}^{M}m\left(\frac1m-\frac1{m+1}\right)
  =\sum_{m=1}^{M}\frac1{m+1}\xrightarrow[M\to\infty]{}+\infty.
$$
これは $G$ の可積分性に反する。よって優収束定理の支配条件は満たされない。
区間の幅が $1/n$ に縮む一方で高さが $n$ に増えるため、積分値は1のまま保たれる。

\clearpage
\handouttitle{第4回\quad 独立性：レポート問題と解答}

\topic{問題4.1\quad 対独立でも相互独立でない例}
\keybox{独立な $X,Y$ がそれぞれ $-1,1$ を確率 $1/2$ でとるとし、$Z=XY$ とおく。
各ペアの同時分布を求め、$X,Y,Z$ が対独立であることを示せ。
次に $\{X=1,Y=1,Z=1\}$ の確率を用い、相互独立ではないことを示せ。}

\examplelabel{解答}
独立性より $\Prob(X=x,Y=y)=1/4$ $(x,y\in\{-1,1\})$ である。
$Z=XY$ なので、三つ組の分布は次の表で尽くされる。
\begin{center}
  \renewcommand{\arraystretch}{1.15}
  \begin{tabular}{@{}cccc@{}}
    \toprule
    $X$ & $Y$ & $Z$ & 確率\\
    \midrule
    $1$ & $1$ & $1$ & $1/4$\\
    $1$ & $-1$ & $-1$ & $1/4$\\
    $-1$ & $1$ & $-1$ & $1/4$\\
    $-1$ & $-1$ & $1$ & $1/4$\\
    \bottomrule
  \end{tabular}
\end{center}
特に $\Prob(Z=1)=\Prob(Z=-1)=1/2$ である。
$x,y,z\in\{-1,1\}$ に対し、$x^2=y^2=1$ を使うと
$$
  \begin{aligned}
    \Prob(X=x,Y=y)&=\frac14,\\
    \Prob(X=x,Z=z)&=\Prob(X=x,Y=xz)=\frac14,\\
    \Prob(Y=y,Z=z)&=\Prob(X=yz,Y=y)=\frac14.
  \end{aligned}
$$
したがって、どのペア $(U,V)\in\{(X,Y),(X,Z),(Y,Z)\}$ についても、
同時分布は $\{-1,1\}^2$ の各点で $1/4$、その外では0である。
各点で $\Prob(U=u,V=v)=\Prob(U=u)\Prob(V=v)$ が成り立つ。
有限個の値しかとらないので、点の確率を足せば任意のボレル集合についても独立性が従う。
よって $X,Y,Z$ は対独立である。しかし
$$
  \Prob(X=1,Y=1,Z=1)=\frac14
  \ne\frac18=\Prob(X=1)\Prob(Y=1)\Prob(Z=1)
$$
であるから、相互独立ではない。

\topic{問題4.2\quad 可測関数による変換と独立性}
\keybox{$X,Y$ を独立な確率変数、$f,g:\R\to\R$ をボレル可測関数とする。
集合の逆像を用いて、$f(X),g(Y)$ が独立であることを示せ。
（$A\in\mathcal B(\mathbb R)$に対し、$f^{-1}(A)\in\mathcal B(\mathbb R)$を用いてよい）}

\examplelabel{解答}
任意の $A,B\in\mathcal B(\R)$ に対して $f^{-1}(A),g^{-1}(B)\in\mathcal B(\R)$ である。
また、$\{f(X)\in A\}=\{X\in f^{-1}(A)\}\in\F$ なので、$f(X)$ は確率変数であり、
$g(Y)$ も同様である。$X,Y$ の独立性から
$$
  \begin{aligned}
    \Prob(f(X)\in A,\ g(Y)\in B)
    &=\Prob(X\in f^{-1}(A),\ Y\in g^{-1}(B))\\
    &=\Prob(X\in f^{-1}(A))\Prob(Y\in g^{-1}(B))\\
    &=\Prob(f(X)\in A)\Prob(g(Y)\in B).
  \end{aligned}
$$
これは任意のボレル集合 $A,B$ で成り立つので、$f(X),g(Y)$ は独立である。

\topic{問題4.3\quad 授業・演習への感想}
\keybox{これまでの授業や演習の感想または改善してほしいところなど。
（ここまででレポート提出、11/3）}

\examplelabel{記入例}
この問題には一意の正解はない。次の例を参考に、自分の経験や考えに合わせて記述する。

「有限の標本空間で確率を計算したあと、連続な標本空間では積分を使う、という流れが
理解の助けになりました。特に、対独立と相互独立の違いは具体例を通して整理できました。
可測性の証明では逆像の書き換えにまだ慣れていないので、授業内でもう一つ例を扱って
いただけると復習しやすいです。」

\label{lastpage}
\end{document}
